\documentclass[11pt]{article}
\usepackage[letterpaper,margin=0.85in]{geometry}\usepackage[T1]{fontenc}\usepackage[utf8]{inputenc}\usepackage{amsmath,amssymb,enumitem,booktabs,array,xcolor,fancyhdr,tabularx}
\setlist{itemsep=2pt,topsep=3pt}\pagestyle{fancy}\fancyhf{}\renewcommand{\headrulewidth}{0.3pt}
\fancyhead[L]{\small 4.7 Part 2 $\cdot$ L'H?pital's Rule: Again, Infinity and When Not to Use It}\fancyhead[R]{\small Teacher key}\fancyfoot[C]{\small\thepage}
\definecolor{ans}{RGB}{28,125,74}\newcommand{\A}[1]{\textcolor{ans}{\textbf{#1}}}
\begin{document}

\section*{4.7 Part 2 $\cdot$ L'H?pital's Rule: Again, Infinity and When Not to Use It --- teacher key}
\textbf{Question:} When is the rule needed twice, how does it work at infinity, and when must I not use it?\par
\subsection*{Self-paced lesson}No clock, no trimming, no skipping, no end mark, no stamp. Every tab is required in order; progress saves on the device so a student can stop and finish the next day.\par
\subsection*{Warm-up review (3 easy items)}\begin{itemize}
\item The derivative of $1-\cos x$ is... \A{$\sin x$.} The derivative of $-\cos x$ is $+\sin x$.
\item $\dfrac{\sin x}{2x}$ at $x=0$ gives... \A{$\dfrac00$.} $\sin 0=0$ and $2\cdot0=0$.
\item $\dfrac{2}{e^x}$ as $x\to\infty$ is... \A{0.} The bottom grows without bound while the top stays 2.
\end{itemize}
\subsection*{Independent lesson}No huddle, no teacher talk, no codes. Students read a short picture tutorial (Concept), watch worked problems move by move (I do), then solve generated problems (You do). Add \texttt{\#key} to the lesson address to open every tab with answers marked.\par
\subsection*{Targets}\begin{itemize}\item I can use L'H?pital's Rule again when the new fraction is still 0/0.
\item I can use the rule for limits as x goes to infinity.
\item I can decide when to stop and when the rule does not apply.\end{itemize}
\subsection*{Predictions}\begin{itemize}
\item As $x$ gets very big, which grows faster, $x^2$ or $e^x$? \A{$e^x$.} An exponential beats any power. The tutorial shows the table.
\item Plugging $x=0$ into $\dfrac{1-\cos x}{x^2}$ gives $\dfrac00$. The limit is... \A{$\dfrac12$.} It takes the rule twice and gives $\dfrac{\cos x}{2}\to\dfrac12$.
\end{itemize}
\section*{Tutorial script, with every question and answer}
\subsection*{Concept tutorial (8 pictures)}
\textit{I can say when to use the rule again, and who wins as x grows.}\par
\begin{enumerate}
\item \textbf{Still 0/0? Go again.} After one use of the rule, plug in again. If it is still $\frac00$, use the rule again. \textit{(Picture: forms. The $\frac00$ row of the table.)}\\ \textit{Question:} For $\dfrac{1-\cos x}{x^2}$ at 0, one use gives $\dfrac{\sin x}{2x}$. At $x=0$ that is... \A{Answer: $\dfrac00$, so use the rule again.} $\sin0=0$ and $2\cdot0=0$, so it is $\frac00$ again.
\item \textbf{The second use.} Use the rule on $\dfrac{\sin x}{2x}$: the top gives $\cos x$ and the bottom gives 2. At $x=0$: $\dfrac{\cos0}{2}=\dfrac12$. \textit{(Picture: forms. Same table row: the form $\frac00$ is highlighted.)}\\ \textit{Question:} Find $\dfrac{\cos x}{2}$ at $x=0$. \A{Answer: 0.5.} $\dfrac{1}{2}=0.5$. This is the limit.
\item \textbf{The form ?/?.} As $x\to\infty$, both $x$ and $e^x$ get huge. That is the form $\frac{\infty}{\infty}$, and the rule applies. \textit{(Picture: growth. Two bars on a log scale: $x$ in blue and $e^x$ in orange. The table lists both values at bigger and bigger $x$. Slide the steps.)}\\ \textit{Question:} In the table, which grows faster: $x$ or $e^x$? \A{Answer: $e^x$.} The orange values explode while the blue values barely grow.
\item \textbf{The ratio shrinks to 0.} At $x=40$ the ratio $x\div e^x$ is almost 0. By the rule: $\dfrac{1}{e^x}\to0$. \textit{(Picture: growth. The last column is $x\div e^x$. Watch it at $x=40$.)}\\ \textit{Question:} $\displaystyle\lim_{x\to\infty}\dfrac{x}{e^x}$ is... \A{Answer: 0.} The bottom grows without bound and the top barely moves, so the ratio goes to 0.
\item \textbf{A bigger power needs more uses.} For $\dfrac{x^2}{e^x}$: $\dfrac{2x}{e^x}$ is still $\frac{\infty}{\infty}$, then $\dfrac{2}{e^x}\to0$. Two uses. \textit{(Picture: growth. Now $x^2$ in blue and $e^x$ in orange.)}\\ \textit{Question:} $\displaystyle\lim_{x\to\infty}\dfrac{x^2}{e^x}$ is... \A{Answer: 0.} After two uses: $\dfrac{2}{e^x}\to0$.
\item \textbf{Logs lose to powers.} The log grows very slowly. $\dfrac{\ln x}{x}$ becomes $\dfrac{1/x}{1}\to0$. \textit{(Picture: growth. $\ln x$ in blue and $x$ in orange.)}\\ \textit{Question:} $\displaystyle\lim_{x\to\infty}\dfrac{\ln x}{x}$ is... \A{Answer: 0.} The log is far slower than $x$, so the ratio goes to 0.
\item \textbf{Two powers.} A higher power wins. $\dfrac{x^2}{x^3}\to0$, and $\dfrac{x^3}{x^2}\to\infty$. \textit{(Picture: growth. $x^2$ in blue and $x^3$ in orange.)}\\ \textit{Question:} $\displaystyle\lim_{x\to\infty}\dfrac{x^2}{x^3}$ is... \A{Answer: 0.} The bottom is the bigger power, so the ratio goes to 0.
\item \textbf{Stop when it is not indeterminate.} Plug in after each use. When you get a plain number, stop: that is the limit, and using the rule again gives a wrong answer. \textit{(Picture: forms. The first row: a plain number.)}\\ \textit{Question:} You use the rule once and get $\dfrac{e^x}{2x+1}$ at $x=0$. What now? \A{Answer: Plug in: $\dfrac11=1$, and stop..} The result is a plain number, so we are done.
\end{enumerate}
\subsection*{I do: the five moves}\begin{enumerate}\item Plug in
\item Use the rule
\item Plug in again
\item Repeat or stop\end{enumerate}
\subsection*{I do (7 steps)}
\textit{I can watch the rule used twice and at infinity, and spot the common slip.}\par
\begin{enumerate}
\item \textbf{Plug in.} \textit{Move 1 $\cdot$ Plug in.} Put $x=0$ into the top and the bottom. \textit{(Picture: forms. The table row for the form $\frac00$ is highlighted.)}\\ Top: $1-\cos0=0$\\ Bottom: $0^2=0$\\ Form: $\dfrac00$\\ \textit{Question:} What is $1-\cos0$? \A{Answer: 0.} $1-1=0$.
\item \textbf{Use the rule.} \textit{Move 2 $\cdot$ Use the rule.} The top becomes $\sin x$ and the bottom becomes $2x$. \textit{(Picture: forms. The $\frac00$ row again: take the top and bottom derivatives separately.)}\\ $\dfrac{\sin x}{2x}$\\ \textit{Question:} Which is the first use of the rule? \A{Answer: $\dfrac{\sin x}{2x}$.} The derivative of $1-\cos x$ is $\sin x$. The derivative of $x^2$ is $2x$.
\item \textbf{Plug in again.} \textit{Move 3 $\cdot$ Plug in again.} Put $x=0$ into $\dfrac{\sin x}{2x}$: $\dfrac00$ again. \textit{(Picture: forms. Same table row: the new fraction is $\frac00$ again.)}\\ $\dfrac{\sin0}{2\cdot0}=\dfrac00$ again\\ \textit{Question:} $\dfrac{\sin x}{2x}$ at $x=0$ is... \A{Answer: $\dfrac00$ again.} Both are 0 again, so use the rule again.
\item \textbf{Repeat the rule.} \textit{Move 4 $\cdot$ Repeat or stop.} Second use: the top gives $\cos x$ and the bottom gives 2, so at $x=0$ it is $\dfrac12$. A plain number: stop. \textit{(Picture: forms. Now a plain number appears.)}\\ $\dfrac{\cos x}{2}\to\dfrac{1}{2}$\\ The limit is $\dfrac12$\\ \textit{Question:} Find $\dfrac{\cos0}{2}$. \A{Answer: 0.5.} $\dfrac12=0.5$.
\item \textbf{Problem 2: infinity.} \textit{Moves 1 and 2 $\cdot$ Plug in, use the rule.} \textbf{Problem:} Find $\displaystyle\lim_{x\to\infty}\dfrac{x^2}{e^x}$. Both the top and the bottom go to $\infty$: the form is $\frac{\infty}{\infty}$. Use the rule. \textit{(Picture: growth. Blue is $x^2$ and orange is $e^x$. Both explode as $x$ grows.)}\\ Form: $\dfrac{\infty}{\infty}$\\ $\dfrac{2x}{e^x}$\\ \textit{Question:} After one use, $\dfrac{2x}{e^x}$ as $x\to\infty$ is... \A{Answer: $\dfrac{\infty}{\infty}$ again.} $2x$ and $e^x$ both grow without bound.
\item \textbf{Again, then stop.} \textit{Moves 3 and 4.} Second use: $\dfrac{2}{e^x}$. The top stays 2 and the bottom grows, so it goes to 0. \textit{(Picture: growth. At $x=40$, $x^2\div e^x$ is nearly 0.)}\\ $\dfrac{2}{e^x}\to0$\\ The limit is $0$\\ \textit{Question:} $\dfrac{2}{e^x}$ as $x\to\infty$ is... \A{Answer: 0.} A fixed top over a huge bottom goes to 0.
\item \textbf{The common slip.} Slip: $\displaystyle\lim_{x\to0}\dfrac{e^x-1}{x^2+x}$ gives $\frac00$, then $\dfrac{e^x}{2x+1}=\dfrac11=1$: STOP, because another use gives the wrong $\dfrac12$. \textit{(Picture: forms. The first row says: a plain number means stop.)}\\ Right: $\dfrac{e^0}{2(0)+1}=1$\\ Wrong: another use gives $\dfrac12$\\ \textit{Common slip.}\\ \textit{Question:} Why is a second use wrong here? \A{Answer: The new fraction is not 0/0 or ?/?, so the rule no longer applies..} Plug in after each use. A plain number means stop.
\end{enumerate}
\subsection*{Big ideas}\begin{itemize}\item After each use, plug in. Still $\frac00$ or $\frac{\infty}{\infty}$: use the rule again.
\item As $x\to\infty$: $\ln x$ is slower than $x^n$, which is slower than $e^x$.
\item A plain number means stop: that number is the limit.
\item A number over 0 is not indeterminate. The rule does not apply.\end{itemize}
\section*{Every level, with answers}
\subsection*{Use the rule more than once}
\noindent\textbf{Level 1.} Find $\displaystyle\lim_{x\to0}\dfrac{1-\cos 2x}{x^2}$. Give it to 2 decimal places if needed.\\ \A{Answer: 2} \quad First use: $\dfrac{2\sin 2x}{2x}$, still $\frac00$. Second use: $\dfrac{4\cos 2x}{2}$, and at $x=0$ that is $\dfrac{4}{2}$.\par\medskip
\noindent\textbf{Level 2.} Find $\displaystyle\lim_{x\to0}\dfrac{e^{x}-1-x}{x^2}$. Give it to 2 decimal places if needed.\\ \A{Answer: 0.5} \quad First use: $\dfrac{e^{x}-1}{2x}$, still $\frac00$. Second use: $\dfrac{1e^{x}}{2}$, and at $x=0$ that is $\dfrac{1}{2}$.\par\medskip
\noindent\textbf{Level 3.} Find $\displaystyle\lim_{x\to0}\dfrac{3x-\sin 3x}{x^3}$. Give it to 2 decimal places if needed.\\ \A{Answer: 4.5} \quad The rule is needed three times: $\dfrac{3-3\cos 3x}{3x^2}$, then $\dfrac{9\sin 3x}{6x}$, then $\dfrac{27\cos 3x}{6}$. At $x=0$ that is $\dfrac{27}{6}$.\par\medskip
\subsection*{Limits at infinity}
\noindent\textbf{Level 1.} Find $\displaystyle\lim_{x\to\infty}\dfrac{x}{e^x}$.\\ \A{(A) $0$} \quad (B) it does not exist \quad (C) $1$ \quad (D) $\infty$\\ \A{Answer: A} \quad $e^x$ beats $x$. By the rule: $\dfrac{1}{e^x}\to0$.\par\medskip
\noindent\textbf{Level 2.} Find $\displaystyle\lim_{x\to\infty}\dfrac{x^3}{e^x}$.\\ (A) $\infty$ \quad \A{(B) $0$} \quad (C) $1$ \quad (D) it does not exist\\ \A{Answer: B} \quad The rule three times gives $\dfrac{6}{e^x}\to0$.\par\medskip
\noindent\textbf{Level 3.} Find $\displaystyle\lim_{x\to\infty}\dfrac{4x^3-x}{2x^3+5}$.\\ (A) $0$ \quad (B) $3$ \quad (C) $\infty$ \quad \A{(D) $2$}\\ \A{Answer: D} \quad The rule three times: $\dfrac{24}{12}=2$.\par\medskip
\subsection*{Again, or stop?}
\noindent\textbf{Level 1.} After one use of the rule you have $\displaystyle\lim_{x\to0}\dfrac{\sin x}{2x}$. What do you do next?\\ \A{(A) Use the rule again.} \quad (B) Stop: the limit is 0 because the top is 0. \quad (C) Stop: it is a number over 0, so the rule does not apply. \quad (D) Stop: plug in, and the limit is the number you get.\\ \A{Answer: A} \quad At $x=0$ the top is 0 and the bottom is 0: $\frac00$ again. So use the rule again.\par\medskip
\noindent\textbf{Level 2.} After one use of the rule you have $\displaystyle\lim_{x\to0}\dfrac{e^x}{2x+1}$. What do you do next?\\ (A) Use the rule again. \quad \A{(B) Stop: plug in, and the limit is the number you get.} \quad (C) Stop: it is a number over 0, so the rule does not apply. \quad (D) Stop: the limit is 0 because the top is 0.\\ \A{Answer: B} \quad At $x=0$ it is $\dfrac11=1$, a plain number. So stop.\par\medskip
\noindent\textbf{Level 3.} After one use of the rule you have $\displaystyle\lim_{x\to\pi}\dfrac{2\cos 2x}{1+\cos x}$. What do you do next?\\ (A) Use the rule again. \quad (B) Stop: the limit is 0 because the top is 0. \quad (C) Stop: plug in, and the limit is the number you get. \quad \A{(D) Stop: it is a number over 0, so the rule does not apply.}\\ \A{Answer: D} \quad At $x=\pi$ the top is 2 and the bottom is 0: a number over 0. The rule does not apply. So stop.\par\medskip
\subsection*{Which one is NOT for the rule?}
\noindent\textbf{Level 1.} Plug in for each limit. Which one is NOT a form where L'H?pital's Rule applies?\\ (A) $\displaystyle\lim_{x\to2}\dfrac{x^3-8}{x-2}$ \quad \A{(B) $\displaystyle\lim_{x\to1}\dfrac{x^2+1}{x-1}$} \quad (C) $\displaystyle\lim_{x\to1}\dfrac{\ln x}{x-1}$ \quad (D) $\displaystyle\lim_{x\to3}\dfrac{x^2-9}{x-3}$\\ \A{Answer: B} \quad Plugging in gives $\dfrac20$, a number over 0. So the rule is not for this limit.\par\medskip
\noindent\textbf{Level 2.} Plug in for each limit. Which one is NOT a form where L'H?pital's Rule applies?\\ (A) $\displaystyle\lim_{x\to0}\dfrac{e^x-1}{x}$ \quad (B) $\displaystyle\lim_{x\to0}\dfrac{\sin x}{x}$ \quad (C) $\displaystyle\lim_{x\to3}\dfrac{x^2-9}{x-3}$ \quad \A{(D) $\displaystyle\lim_{x\to2}\dfrac{x+4}{x+1}$}\\ \A{Answer: D} \quad Plugging in gives $\dfrac63=2$, a plain number. So the rule is not for this limit.\par\medskip
\noindent\textbf{Level 3.} Plug in for each limit. Which one is NOT a form where L'H?pital's Rule applies?\\ (A) $\displaystyle\lim_{x\to1}\dfrac{\ln x}{x-1}$ \quad (B) $\displaystyle\lim_{x\to2}\dfrac{x^3-8}{x-2}$ \quad \A{(C) $\displaystyle\lim_{x\to2}\dfrac{x^2-4}{x+3}$} \quad (D) $\displaystyle\lim_{x\to0}\dfrac{e^x-1}{x}$\\ \A{Answer: C} \quad Plugging in gives $\dfrac05=0$, a plain number. So the rule is not for this limit.\par\medskip
\section*{Common misconceptions}\begin{tabularx}{\textwidth}{@{}X X@{}}\toprule Students say & Say instead\\\midrule
Use the rule again even when the form is a plain number. & Plug in after each use. A plain number means stop.\\
?/? means the limit is 1. & ?/? is indeterminate. The growth rates decide: 0, ? or a number.\\
A number over 0 can use the rule. & It is not 0/0 or ?/?. The rule does not apply.\\
The quotient rule is part of the rule. & Differentiate the top and the bottom separately.\\
x? beats e? as x grows. & e? beats every power of x. Every power of x beats ln x.\\
\bottomrule\end{tabularx}
\end{document}